\documentclass[a4paper, 14pt]{extarticle} % 14pt вместо 12pt
\usepackage[T2A]{fontenc}
\usepackage[utf8]{inputenc}
\usepackage[russian]{babel}
\usepackage{newtxtext,newtxmath} % Times New Roman шрифт
\usepackage{geometry}
\usepackage{tabularx}
\usepackage{multirow}
\usepackage{booktabs}
\usepackage{amsmath}
\usepackage{graphicx}
\usepackage{float}
\usepackage{listings}
\usepackage{xcolor}
\usepackage{setspace} % Для межстрочного интервала
\usepackage{indentfirst} % Отступ первой строки абзаца

% Настройка страницы по ГОСТ
\geometry{
    left=30mm,    % 3 см
    right=15mm,   % 1,5 см
    top=20mm,     % 2 см
    bottom=20mm,  % 2 см
}

% Межстрочный интервал 1,5
\onehalfspacing

% Отступ первой строки абзаца 1,25 см
\setlength{\parindent}{1.25cm}

% Настройка стиля листингов
\lstset{
    language=C++,
    basicstyle=\ttfamily\small, % немного меньше основного текста
    keywordstyle=\color{blue}\bfseries,
    commentstyle=\color{green!60!black},
    stringstyle=\color{red},
    numbers=left,
    numberstyle=\tiny\color{gray},
    stepnumber=1,
    numbersep=5pt,
    tabsize=4,
    showstringspaces=false,
    breaklines=true,
    frame=single,
    backgroundcolor=\color{gray!5},
    rulecolor=\color{gray!30},
    captionpos=b,
    aboveskip=10pt,
    belowskip=10pt
}

% Настройка заголовков
\usepackage{titlesec}

% Формат заголовков разделов (section)
\titleformat{\section}
    {\normalfont\Large\bfseries}  % размер: Large, шрифт: жирный
    {\thesection}                  % номер раздела
    {1em}                          % расстояние между номером и названием
    {}                             % код после заголовка

% Формат подразделов (subsection)
\titleformat{\subsection}
    {\normalfont\large\bfseries}   % размер: large, шрифт: жирный
    {\thesubsection}
    {1em}
    {}

% Формат под-подразделов (subsubsection)
\titleformat{\subsubsection}
    {\normalfont\normalsize\bfseries} % размер: normal, шрифт: жирный
    {\thesubsubsection}
    {1em}
    {}

% Установка отступов до и после заголовков
\titlespacing*{\section}{0pt}{12pt}{6pt}      % {влево}{до}{после}
\titlespacing*{\subsection}{0pt}{10pt}{4pt}
\titlespacing*{\subsubsection}{0pt}{8pt}{2pt}

% Убираем отступы в списках
\usepackage{enumitem}
\setlist{nosep, leftmargin=1.25cm}

% Для не нумерованных заголовков тоже
\titleformat{name=\section,numberless}
    {\normalfont\Large\bfseries}
    {}
    {0pt}
    {}

\titleformat{name=\subsection,numberless}
    {\normalfont\large\bfseries}
    {}
    {0pt}
    {}

% Настройка гиперссылок (рекомендуется подключать последним)
\usepackage{hyperref}
\hypersetup{
    colorlinks=true,
    linkcolor=blue,
    filecolor=magenta,      
    urlcolor=cyan,
    pdftitle={Отчет по лабораторной работе №2},
    pdfauthor={Храмов В.Д.}
}

%%%%%%%%%%%%%%%%%%%%%%%%%%%%%%%%%%%%%%%%%%%%%%%%%%%%%%%%%%%%%%%

\begin{document}

    \begin{titlepage}
        \thispagestyle{empty}
        
        \begin{center}
        
            { \large МИНИСТЕРСТВО НАУКИ И ВЫСШЕГО ОБРАЗОВАНИЯ \\ 
                РОССИЙСКОЙ ФЕДЕРАЦИИ }
             
            { «Санкт-Петербургский политехнический университет Петра Великого» }
            
            { Институт компьютерных наук и кибербезопасности \\ 
                Высшая школа технологий искусственного интеллекта \\
                Направление: 02.03.01. «Математика и компьютерные науки» \\ }        
              
              \vspace{3cm}
            { \large Техническое задание к курсовой работе \\ 
              По дисциплине: «Функциональное программирование» \\ 
              На тему: «Учет рабочего времени сотрудников с расчетом зарплаты»}
    
        \end{center}
    
        \vspace{3cm}

        \begin{center}
            % Студент
            \noindent
            \begin{minipage}[t]{\textwidth}
                \raggedright Студент  \\ группы 5130201/40002 \hfill 
                \underline{\hspace{3cm}} Храмов В.Д.
            \end{minipage}
            
            \vspace{1cm}
            
            % Преподаватель
            \noindent
            \begin{minipage}[t]{\textwidth}
                \raggedright к. т. н. \hfill 
                \underline{\hspace{3cm}} Моторин Д.Е.
            \end{minipage}
        \end{center}
        
        \vspace{1cm}
        
        % Дата справа
        \begin{flushright}
            \large
            «\underline{\hspace{1cm}}» \underline{\hspace{3cm}} 2026 г.
        \end{flushright}
        
        \vspace{0.5cm}
        
        \vfill
        
        \begin{center}
            Санкт-Петербург \\
            2026
        \end{center}   
            
    \end{titlepage}
    
    \pagenumbering{arabic} 
    \setcounter{page}{2}  
%%%%%%%%%%%%%%%%%%%%%%%%%%%%%%%%%%%%%%%%%%%%%%%%%%%%%%%%%%%%%%%
    \newpage
    \tableofcontents
%%%%%%%%%%%%%%%%%%%%%%%%%%%%%%%%%%%%%%%%%%%%%%%%%%%%%%%%%%%%%%%
    \newpage
    \section{Основание для выполнения работы}
        Работа выполняется на основании учебного плана направления подготовки 02.03.01 «Математика и компьютерные науки» и рабочей программы дисциплины «Функциональное программирование».
%%%%%%%%%%%%%%%%%%%%%%%%%%%%%%%%%%%%%%%%%%%%%%%%%%%%%%%%%%%%%%%
    \newpage
    \section{Исполнитель и соисполнители работы}   
        Исполнитель: студент 2 курса, группы 5130201/40002, Храмов Владислав Дмитриевич
        
        Руководитель работы: кандидат технических наук, Моторин Дмитрий Евгеньевич
%%%%%%%%%%%%%%%%%%%%%%%%%%%%%%%%%%%%%%%%%%%%%%%%%%%%%%%%%%%%%%%
    \newpage
    \section{Цели, задачи и исходные данные}
        \subsection{Цель работы}
            Целью работы является разработка программы на языке Haskell для автоматизации учёта рабочего времени сотрудников и расчёта заработной платы с возможностью формирования отчётов и логирования действий пользователя.
        \subsection{Задачи работы}
            В процессе выполнения работы требуется решить следующие задачи:
            \begin{enumerate}
                \item Провести анализ предметной области (учёт рабочего времени, системы оплаты труда).
                \item Разработать структуры данных для представления информации о сотрудниках, учете рабочего времени, ставках оплаты.
                \item Реализовать модуль загрузки данных из файлов формата CSV и сохранения результатов в файлы.
                \item Разработать математический аппарат для расчёта:
                    \begin{itemize}
                        \item отработанных часов (обычных, ночных, сверхурочных);
                        \item заработной платы с учётом различных систем оплаты;
                        \item налогов и отчислений (НДФЛ, страховые взносы).
                    \end{itemize}
                \item Создать консольный интерфейс для взаимодействия с пользователем (меню, ввод параметров, вывод результатов).
                \item Реализовать систему логирования действий пользователя с временными метками.
                \item Обеспечить обработку ошибочных ситуаций (отсутствие файлов, некорректные данные, несоответствие форматов).
                \item Разработать набор тестов (с использованием QuickCheck) для проверки корректности ключевых функций.
                \item Подготовить отчёт по курсовой работе.
            \end{enumerate}
        \subsection{Исходные данные}
            Исходными данными для проведения работы являются:
                \begin{itemize}
                    \item Табели учёта рабочего времени в формате CSV (поля: табельный номер, дата, время прихода, время ухода).
                    \item Справочник сотрудников в формате CSV (поля: табельный номер, ФИО, должность, система оплаты, ставка).
                    \item Производственный календарь (список праздничных дней) в текстовом формате.
                    \item Конфигурационный файл с параметрами расчёта (норма часов, коэффициенты оплаты, ставки налогов).
                \end{itemize}
%%%%%%%%%%%%%%%%%%%%%%%%%%%%%%%%%%%%%%%%%%%%%%%%%%%%%%%%%%%%%%%
    \newpage
    \section{Основное содержание работы}
        \subsection{Разработка проекта}
            В ходе выполнения работы разрабатывается:
            \begin{itemize}
                \item Программный комплекс «Учёт рабочего времени и расчёт зарплаты» (далее — ПК), реализованный на языке Haskell.
                \item Исходные тексты программы, организованные в логические модули (например, Data.Types, Data.Parser, Logic.Salary, Logic.Overtime, IO.Console, IO.Logger, Main).
                \item Тестовые наборы данных (примеры файлов) для проверки функциональности.
                \item Документация по установке, настройке и использованию.
            \end{itemize}







        \subsection{Отчетная документация}
            \begin{itemize}
                \item Техническое задание
                \item Диаграмма типов и модулей
                \item Отчет
            \end{itemize}
%%%%%%%%%%%%%%%%%%%%%%%%%%%%%%%%%%%%%%%%%%%%%%%%%%%%%%%%%%%%%%%
    \newpage
    \section{Основные требования к выполнению работы}
        \subsection{Назначение разрабатываемого программного комплекса}
            Программный комплекс предназначен для автоматизации учёта рабочего времени сотрудников предприятия и расчёта заработной платы. Программа может использоваться в отделах кадров и бухгалтерии малых предприятий, а также в учебных целях для изучения функционального программирования.
        \subsection{Состав разрабатываемого программного комплекса}
            ПК состоит из следующих модулей:
                \begin{itemize}
                    \item Модуль загрузки данных – чтение CSV-файлов, парсинг дат/времени, преобразование во внутренние структуры.
                    \item Модуль расчёта времени – вычисление отработанных часов, выделение ночных, сверхурочных, праздничных часов.
                    \item Модуль расчёта зарплаты – вычисление заработной платы по каждому сотруднику за выбранный период с учётом системы оплаты и налогов.
                    \item Модуль консольного интерфейса – организация диалога с пользователем, отображение меню, ввод данных, вывод таблиц.
                    \item Модуль логирования – запись действий пользователя в лог-файл.
                    \item Модуль обработки ошибок – функции для безопасной обработки исключительных ситуаций.
                    \item Модуль тестирования – набор модульных тестов и свойств для QuickCheck
                \end{itemize}
            \subsection{Технические характеристики}
                Программное обеспечение должно обеспечить следующие технические характеристики по назначению:
                \begin{itemize}
                    \item Время загрузки файла объёмом до 10 000 записей – не более 2 секунд.
                    \item Время расчёта зарплаты для 100 сотрудников за месяц – не более 1 секунды.
                    \item Объём оперативной памяти, занимаемой программой при работе с типовыми данными, – не более 100 МБ.
                    \item Корректная работа под управлением операционных систем семейства Linux, macOS, Windows (при наличии компилятора GHC).
                \end{itemize}

            \subsection{Функциональные требования}
                Программа должна обеспечивать выполнение следующих функций:
                \begin{enumerate}
                    \item Загрузка данных:
                        \begin{itemize}
                            \item Загрузка справочника сотрудников из CSV-файла.
                            \item Загрузка табелей рабочего времени из CSV-файла (можно несколько файлов за разные периоды).
                            \item Загрузка производственного календаря из текстового файла.
                            \item Загрузка конфигурации из JSON/YAML-файла.
                        \end{itemize}
                    \item Расчёт:
                        \begin{itemize}
                            \item Расчёт количества отработанных часов за период (день, неделя, месяц) для каждого сотрудника.
                            \item Выделение часов, отработанных в ночное время (если задано в конфигурации).
                            \item Выделение часов, отработанных в праздничные дни.
                            \item Расчёт сверхурочных часов (сверх нормы за день/неделю/месяц).
                            \item Расчёт заработной платы:
                                \begin{itemize}
                                    \item для почасовой оплаты: часы * ставка;
                                    \item для окладной: оклад * (отработанные часы / норма часов) с учётом переработок;
                                    \item для сдельной: количество продукции * расценка (если реализовано).
                                    \item Расчёт НДФЛ и страховых взносов по заданным ставкам.
                                \end{itemize}
                        \end{itemize}
                    \item Вывод результатов:
                        \begin{itemize}
                            \item Отображение сводной таблицы по сотрудникам (ФИО, отработанные часы, зарплата, налоги, сумма к выдаче) на экране.
                            \item Сохранение отчёта в текстовый файл.
                            \item Возможность фильтрации по подразделению/должности (при наличии соответствующих данных).
                        \end{itemize}
                    \item Логирование:
                        \begin{itemize}
                            \item Запись в лог-файл всех значимых действий: запуск программы, загрузка файлов, выполнение расчётов, сохранение отчётов, ошибки.
                            \item Каждая запись должна содержать дату и время события.
                        \end{itemize}
                    \item Взаимодействие с пользователем:
                            \begin{itemize}
                                \item Интерактивное меню с пунктами:
                                    \begin{itemize}
                                        \item Загрузить данные
                                        \item Выбрать период расчёта
                                        \item Выполнить расчёт зарплаты
                                        \item Показать отчёт
                                        \item Сохранить отчёт в файл
                                        \item Просмотреть лог
                                        \item Выход
                                    \end{itemize}
                                \item Запрос необходимых параметров (например, начальная и конечная дата периода).
                                \item Вывод сообщений об ошибках с понятными подсказками.
                            \end{itemize}
                    \item Обработка ошибок ввода и данных:
                        \begin{itemize}
                            \item Проверка существования файлов при загрузке; если файл не найден – вывод сообщения и возврат в меню.
                            \item Проверка корректности формата данных в CSV (правильное количество полей, допустимые значения дат/времени, числа). Некорректные строки должны пропускаться с занесением информации в лог.
                            \item Обработка ситуаций, когда загруженные данные отсутствуют или недостаточны для расчёта (например, нет табелей за выбранный период) – вывод предупреждения.
                            \item Использование Maybe и Either для безопасных вычислений (например, деление на ноль при расчёте нормы часов).
                            \item Защита от переполнения при работе с большими числами (использование типов с фиксированной точностью или Integer для денежных сумм).
                        \end{itemize}
                \end{enumerate}
%%%%%%%%%%%%%%%%%%%%%%%%%%%%%%%%%%%%%%%%%%%%%%%%%%%%%%%%%%%%%%%
    \newpage
    \section{Этапы и сроки выполнения работы}

%%%%%%%%%%%%%%%%%%%%%%%%%%%%%%%%%%%%%%%%%%%%%%%%%%%%%%%%%%%%%%%
    \newpage
    \section{Результаты работы}

%%%%%%%%%%%%%%%%%%%%%%%%%%%%%%%%%%%%%%%%%%%%%%%%%%%%%%%%%%%%%%%
\end{document}
%%%%%%%%%%%%%%%%%%%%%%%%%%%%%%%%%%%%%%%%%%%%%%%%%%%%%%%%%%%%%%%
